\documentclass[12pt]{article}

%%% These are some packages that are useful
\usepackage{amsfonts, lipsum}
\usepackage[scr]{rsfso}
\usepackage{amsmath,amssymb, amscd,amsbsy, bbm, amsthm, enumerate}
\usepackage[top=1in, bottom=1in, left=.45in, right=.45in]{geometry}
\usepackage[unicode]{hyperref}
\usepackage{tikz, xcolor, fancyhdr}
\usepackage{graphicx, caption, subcaption}

%%% Page formatting
%\setlength{\headsep}{30pt}
\setlength{\parindent}{0pt}
\setlength{\textheight}{9in}



%%% These define our question environment and help number things correctly
\theoremstyle{definition}
\newtheorem{thm}{Theorem}
\newtheorem{question}[thm]{Question}
\newtheorem{questionS}[thm]{***Question}
\newtheorem{prop}[thm]{Proposition}
\newtheorem{lem}[thm]{Lemma}
\newtheorem{cor}[thm]{Corollary}
\newtheorem{defn}[thm]{Definition}
\newtheorem{rem}[thm]{Remark}
\newtheorem{ex}[thm]{Example}

%%% These are some shortcuts that are handy
\def\real{{\mathbb R}}
\def\Natural{\mathbb{N}}
\def\dx{\textnormal{dx}}
\def\dy{\textnormal{dy}}
\def\dz{\textnormal{dz}}
\def\dt{\textnormal{dt}}
\def\Re{\textnormal{Re}}
\def\Im{\textnormal{Im}}
\def\exp{\textnormal{exp}}
\def\interior{\textnormal{interior}}
\def\al{\alpha}
\def\del{\delta}
\def\Del{\Delta}
\def\gam{\gamma}
\def\Gam{\Gamma}
\def\Om{\Omega}
\def\ep{\varepsilon}
\def\lam{\lambda}
\def\rational{{\mathbb Q}}
\def\integer{{\mathbb Z}}
\def\Q{{\mathbb Q}}
\def\Z{{\mathbb Z}}
\def\N{{\mathbb N}}
\def\R{{\mathbb R}}
\newcommand{\F}{\mathbb{F}}
\def\grad{\nabla}
\def\C{\mathcal C}
\def\P{\mathbb{P}}
\def\T{\mathcal T}
\def\I{\mathcal I}
\newcommand{\abs}[1]{\left| #1 \right|}
\newcommand{\inner}[1]{\langle #1 \rangle}
\newcommand{\norm}[1]{\left\lVert#1\right\rVert}
\newcommand{\spanvect}{\textnormal{span}}
\newcommand{\union}{\cup}
\newcommand{\Union}{\bigcup}
\def\intersect{\cap}
\def\Intersect{\bigcap}
\renewcommand{\neg}{\mathord{\sim}}
\newcommand{\ds}{\displaystyle}

\DeclareMathOperator{\dg}{deg}





%%%%%%%%%%%%%%%%%%%%%%%%%%%%%%%%%
%%%%%%%%%%%%%%%%%%%%%%%%%%%%%%%%%
\newcommand{\name}{Duff Baker-Jarvis}
%%%%%%%%%%%%%%%%%%%%%%%%%%%%%%%%%
%%%%%%%%%%%%%%%%%%%%%%%%%%%%%%%%%
%%% Document Starts now
\begin{document}
\title{Spectral Theory of Graphs}
\date{\today}
\author{Dr. Duff Baker-Jarvis}
\maketitle

\section*{Primary References}
\begin{enumerate}
    \item \textit{Spectral and Algebraic Graph Theory} by Daniel Spielman
    \item \textit{Spectral Graph Theory and Related Topics} Lectures by Matthew Hirn
\end{enumerate}


\section*{Basic Definitions}
\begin{defn}
    A \textit{Graph} is a pair $ (E,V) $ such that $ V $ is a set of vertices and $ E\subseteq \binom{V}{2}$. Note that this definition disallows loops and multiple edges.
    A \textit{Weighted Graph} is a triple $ (E,V, f) $ where $ (E,V) $ is a graph and $ f:E\to \F $. We will primarily use weights in $ \R $.
\end{defn}

\begin{defn}
The adjacency matrix of a finite graph is 
$$ \mathbf{M}_G(a,b)= \begin{cases}
    1, & (a,b)\in E\\
    0, & (a,b)\not\in E
\end{cases}
 $$
 We can also modify this by changing $ 1\to w(a,b) $ for a weighted graph.
\end{defn}

\begin{defn}
For a vertex $ a $ in a graph $ (E,V) $, the neighborhood of $ a $ is the set of all vertexes connected to $ a $ by an edge.
$$ N_G(a)=\left\{b\in V\; :\; (a,b)\in E\right\}. $$ 
The degree of a vertex $ a $ is 
$$ \dg(a) = |N_G(a)|. $$ 
If it is a weighted graph, we add up all the weights of the incident edges.
\end{defn}

\begin{defn}
The degree vector $ \mathbf{d} $ is given by
$$ \mathbf{M} \mathbf{1}.$$ 
\end{defn}


\begin{defn}
The degree matrix of a graph is obtained by placing the degrees in a diagonal matrix.
$$ \mathbf{D}_G(a,a) = \mathbf{d}(a) $$ 
and 0 on all off diagonals.
\end{defn}

\begin{defn}
The random walk matrix of a graph is 
$$ \mathbf{W}_G = \mathbf{M}_G \mathbf{D}_G^{-1} $$ 
\end{defn}

\section*{The Graph Laplacian}
The normal laplacian $\ds \Delta f = \sum_{k=1}^{n} \frac{\partial^2 f}{\partial x_k}  $. In one dimension, this is simply $ \Delta f = f'' $.
\begin{lem}
Note that 
$$ f'(c)= \lim_{h\to 0} \frac{f(c+h)-f(c-h)}{2h}. $$ 
\end{lem}

Proof Sketch: Add a middle term $ f(c)-f(c) $ and obtain two copies of the definition of derivative (one with $ h $ replaced with $ -h $.

\begin{lem}
Observe that 
$$ f''(c)=\Delta f(c)= \lim_{h\to 0} \frac{f(c+h)+f(c-h)-2f(c)}{h^2} $$ 
\end{lem}

Proof Sketch: Use L'Hospital's rule to turn it into previous lemma.



\end{document} 
